\documentclass[11pt,a4paper]{article}
\usepackage[margin=1.2in]{geometry}
\usepackage{xcolor}
\usepackage{hyperref}
\usepackage{enumitem}
\usepackage{titlesec}
\usepackage{tabularx}
\usepackage{booktabs}
\usepackage{parskip}

% Color definitions based on the flyer
\definecolor{primary}{HTML}{0A3044}
\definecolor{secondary}{HTML}{D4AF37}
\definecolor{accent}{HTML}{145A7A}

% Hyperlink setup
\hypersetup{
    colorlinks=true,
    linkcolor=primary,
    filecolor=magenta,      
    urlcolor=accent,
}

% Section styling
\titleformat{\section}
  {\Large\bfseries\color{primary}}
  {}{0em}{}
  [\vspace{-0.3em}\rule{\textwidth}{1pt}]

\titleformat{\subsection}
  {\large\bfseries\color{accent}}
  {}{0em}{}

\begin{document}

\begin{center}
    {\Large \textbf{\color{primary} MALAVIYA MISSION TEACHER TRAINING CENTRE}}\\[0.3em]
    {\LARGE \textbf{\color{primary} KANNUR UNIVERSITY}}\\[0.3em]
    {\small Thavakkara, Civil Station. P.O., Kannur District, Kerala - 670 002}\\[1.5em]
    
    {\Large \textbf{\color{secondary} ONLINE REFRESHER COURSE ON}}\\[0.5em]
    {\Huge \textbf{\color{primary} Mathematical Foundations}}\\[0.3em]
    {\Large \textbf{\color{primary} of Contemporary Scientific Research}}\\[1em]
    {\large \textbf{15-July-2026 to 28-July-2026}}\\[1.5em]
\end{center}

\section*{Programme Philosophy \& Objectives}
The central philosophy of this refresher course is to present mathematics as a foundational language of modern science and technology, viewed from the perspective of a broad scientific audience rather than that of specialists alone. 

The primary objectives include:
\begin{itemize}[label=\textcolor{secondary}{$\bullet$}]
    \item Exposing participants to important mathematical ideas, themes, and developments that significantly influence contemporary scientific research.
    \item Ensuring concepts remain accessible to faculty members from diverse academic backgrounds.
    \item Emphasizing intuition, motivation, practical applications, historical development, and cross-disciplinary connections with science and technology over strictly technical details and rigorous proofs.
    \item Helping participants appreciate the overarching beauty, power, and relevance of mathematics across various disciplines.
\end{itemize}

\section*{Target Audience}
Faculty members from the following streams are eligible and encouraged to apply:
\begin{itemize}[label=\textcolor{secondary}{$\bullet$}, itemsep=0pt]
    \item Mathematics
    \item Statistics
    \item Physics
    \item Chemistry
    \item Computer Science
    \item Other allied Science and Technology disciplines
\end{itemize}

\section*{Course Structure \& Key Features}
\begin{itemize}[label=\textcolor{secondary}{$\bullet$}, itemsep=0pt]
    \item \textbf{Total Duration:} 12 working days
    \item \textbf{Daily Schedule:} 4 daily sessions
    \item \textbf{Session Length:} 90 minutes per session 
    \item \textbf{Career Advancement:} Considered for CAS
    \item \textbf{Registration Fee:} None (No Registration Fee)
    \item \textbf{Mode of Conduct:} Online
\end{itemize}

\section*{Registration \& Contact Details}

\begin{minipage}[t]{0.5\textwidth}
    \textbf{\color{accent} Registration Portals:}\\[0.3em]
    \url{https://mmc.ugc.ac.in/}\\
    \url{https://mmttc.kannuruniversity.ac.in/}
\end{minipage}
\begin{minipage}[t]{0.5\textwidth}
    \textbf{\color{accent} Programme Director:}\\[0.3em]
    MMTTC Kannur University\\
    Phone: 0497-2700368, 9446353952\\
    Email: \href{mailto:ugchrdc@kannuruniv.ac.in}{ugchrdc@kannuruniv.ac.in}
\end{minipage}

\vspace{1.5em}

\section*{Confirmed Speakers \& Resource Persons}
\textit{Note: Lecture topics and exact schedules for each speaker will be updated continuously as they are finalized. Expected format for speakers is one or two 90-minute expository sessions.}

\vspace{0.5em}

\begin{tabularx}{\textwidth}{@{}l X c@{}}
\toprule
\textbf{Speaker} & \textbf{Title / Topic} & \textbf{Sessions} \\
\midrule
1. Pranav & \textit{To be updated} & 2 \\
2. Renjith & \textit{To be updated} & 2 \\
3. Nikhil & \textit{To be updated} & 1 \\
4. Muruganandam & \textit{To be updated} & 2 \\
5. Kaushal Verma & \textit{To be updated} & 2 \\
6. Krishna Hanumanthu & \textit{To be updated} & 1 or 2 \\
7. Jaydeb Sarkar & \textit{To be updated} & 2 \\
8. Prasuna & \textit{To be updated} & 1 or 2 \\
9. Ajith Parameswaran & \textit{To be updated} & 2 \\
10. Nandakumaran & \textit{To be updated} & 2 \\
11. KNR & \textit{To be updated} & 2 \\
12. Lakshmipriya & \textit{To be updated} & 2 \\
13. Tulsiram & \textit{To be updated} & 2 \\
14. Govind & \textit{To be updated} & 1 or 2 \\
15. Jishnu & \textit{To be updated} & 1 \\
16. Abhishek & \textit{To be updated} & 1 \\
\bottomrule
\end{tabularx}

\vspace{1.5em}

\section*{Programme Schedule}
\begin{center}
    \textit{[Detailed day-by-day session schedule to be prepared and inserted here.]}
\end{center}

\end{document}
